\documentclass[french]{article}

\usepackage[utf8]{inputenc}
\usepackage[T1]{fontenc}
\usepackage{lmodern}
\usepackage{geometry}
\usepackage{graphicx}
\usepackage{babel}
\usepackage{amsmath}
\usepackage{amsthm}
\usepackage{amssymb}
\usepackage{hyperref}
\usepackage{stmaryrd}
\usepackage{enumitem}
\usepackage[most]{tcolorbox}
\usepackage{dsfont}
\usepackage{multirow}
\usepackage{etoc}
\usepackage{titlesec}
\usepackage{esint}
\usepackage{cancel}
\usepackage{mathdots}
\usepackage{indentfirst}
\usepackage{lastpage}
\usepackage{fancyhdr}

\pagestyle{fancy}
\fancyhf{}
\renewcommand{\headrulewidth}{0pt}
\renewcommand{\footrulewidth}{0pt}
\fancyhead[L,R]{}
\fancyfoot[R]{\thepage \ /\ \pageref{LastPage}}

\fancypagestyle{plain}{\fancyhead{}\renewcommand{\headrule}{}}

\setlength{\parindent}{0pt}

\setlist[itemize]{label=$\bullet$}
\setlist{before=\leavevmode, leftmargin=*}

\renewcommand{\P}{\mathbb{P}}
\newcommand{\N}{\mathbb{N}}
\newcommand{\Z}{\mathbb{Z}}
\newcommand{\Q}{\mathbb{Q}}
\newcommand{\R}{\mathbb{R}}
\newcommand{\C}{\mathbb{C}}
\newcommand{\K}{\mathbb{K}}
\renewcommand{\L}{\mathbb{L}}
\newcommand{\U}{\mathbb{U}}
\newcommand{\st}{^*}
\newcommand{\tq}{\middle|}
\newcommand{\inv}{^{-1}}
\newcommand{\E}{\mathbb{E}}
\newcommand{\F}{\mathbb{F}}
\newcommand{\V}{\mathbb{V}}
\renewcommand{\ker}{\text{Ker}}
\newcommand{\im}{\text{Im}}
\newcommand{\card}{\text{Card}}
\newcommand{\Tr}{\text{Tr}}
\newcommand{\Vect}{\text{Vect}}
\newcommand{\rg}{\text{rg}}
\newcommand{\vertiii}[1]{{\left\vert\kern-0.25ex\left\vert\kern-0.25ex\left\vert #1 \right\vert\kern-0.25ex\right\vert\kern-0.25ex\right\vert}}
\renewcommand{\o}{\text{o}}
\renewcommand{\O}{\text{O}}
\renewcommand{\d}{\text{d}}
\newcommand{\Id}{\text{Id}}
\newcommand{\Sp}{\text{Sp}}


\theoremstyle{plain}
\newtheorem*{ind}{Indication}

\theoremstyle{definition}

\newtheorem*{ex}{Exemple}
\newtheorem*{rmq}{Remarque}
\newtheorem*{notation}{Notation}
\newtheorem*{rappel}{Rappel}
\newtheorem*{terminologie}{Terminologie}
\newtheorem*{attention}{Attention}
\newtheorem*{conséquence}{Conséquence}

\theoremstyle{remark}
\newtheorem*{app}{Application}

\newtcolorbox[auto counter]{dfn}[1][]{
	enhanced,
	colback=white,
	colframe=black,
	sharp corners,
	boxrule=0.8pt,
	fonttitle=\bfseries,
	colbacktitle=white,
	coltitle=black,
	attach boxed title to top left={yshift=-1mm,xshift=-0.5mm},
	boxed title style={size=minimal, right=2pt, sharp corners, colframe=white},
	title={Définition \thetcbcounter \ifstrempty{#1}{}{ (#1)}},
	before skip=0.5em,
	after skip=0.5em
}

\newtcolorbox[auto counter]{exo}[1][]{
	enhanced,
	arc=1mm,
	colback=white,
	colframe=black,
	coltitle=black,
	fonttitle=\bfseries,
	boxrule=0.8pt,
	shadow={1.2pt}{-1.2pt}{0pt}{black!50},
	attach title to upper={\ },
	title={Exercice \thetcbcounter\ifblank{#1}{}{ #1}},
	before skip=0.5em,
	after skip=0.5em,
	left=2mm, right=2mm, top=1mm, bottom=1mm
}

\newtcolorbox{met}[1][]{
	enhanced,
	arc=1mm,
	colback=white,
	colframe=black,
	coltitle=black,
	fonttitle=\bfseries,
	boxrule=0.8pt,
	shadow={1.2pt}{-1.2pt}{0pt}{black!50},
	attach title to upper={\ },
	title={Point méthode \ifblank{#1}{}{#1}},
	after title={\par\nobreak\medskip\noindent},
	before skip=1em,
	after skip=1em,
	left=2mm, right=2mm, top=1mm, bottom=1mm
}

\newcounter{thmcount}

\newtcolorbox[use counter=thmcount]{thm}[1][]{
	enhanced,
	colback=white,
	colframe=black,
	sharp corners,
	left skip=0mm,
	boxrule=1.3pt,
	fonttitle=\bfseries,
	colbacktitle=white,
	coltitle=black,
	attach boxed title to top left={yshift=-1mm,xshift=-0.5mm},
	boxed title style={size=minimal, right=2pt, sharp corners, colframe=white},
	title={Théorème \thethmcount \ifstrempty{#1}{}{ (#1)}},
	before skip=0.5em,
	after skip=0.5em
}

\newtcolorbox[use counter=thmcount]{prop}[1][]{
	enhanced,
	colback=white,
	colframe=black,
	sharp corners,
	boxrule=0.8pt,
	fonttitle=\bfseries,
	colbacktitle=white,
	coltitle=black,
	attach boxed title to top left={yshift=-1mm,xshift=-0.5mm},
	boxed title style={size=minimal, right=2pt, sharp corners, colframe=white},
	title={Proposition \thethmcount \ifstrempty{#1}{}{ (#1)}},
	before skip=0.5em,
	after skip=0.5em
}

\newtcolorbox[use counter=thmcount]{cor}[1][]{
	enhanced,
	colback=white,
	colframe=black,
	sharp corners,
	boxrule=0.8pt,
	fonttitle=\bfseries,
	colbacktitle=white,
	coltitle=black,
	attach boxed title to top left={yshift=-1mm,xshift=-0.5mm},
	boxed title style={size=minimal, right=2pt, sharp corners, colframe=white},
	title={Corollaire \thethmcount \ifstrempty{#1}{}{ (#1)}},
	before skip=0.5em,
	after skip=0.5em
}

\newtcolorbox[use counter=thmcount]{lem}[1][]{
	enhanced,
	colback=white,
	colframe=black,
	sharp corners,
	boxrule=0.5pt,
	fonttitle=\bfseries,
	colbacktitle=white,
	coltitle=black,
	attach boxed title to top left={yshift=-1mm,xshift=-0.5mm},
	boxed title style={size=minimal, right=2pt, sharp corners, colframe=white},
	title={Lemme \thethmcount \ifstrempty{#1}{}{ (#1)}},
	before skip=0.5em,
	after skip=0.5em
}

\addto\captionsfrench{\renewcommand*{\contentsname}{Réduction des endomorphismes}}

\title{Réduction des endomorphismes}
\date{}

\begin{document}
	\tableofcontents
	
	\newpage
	\maketitle
	
	Dans ce chapitre, $E$ est un espace vectoriel de dimension quelconque sur un sur-corps $\K$ de $\Q$. On se limite dans la pratique au cas où $\K$ est $\R$ ou $\C$.\\
	On étudie, dans ca chapitre, des endomorphismes de $E$ et, en particulier, lorsque $E$ est de dimension finie, on cherche, si c'est possible, une base de $E$ dans laquelle un tel endomorphisme possède une matrice simple : diagonale, voire triangulaire. \\
	Cela revient, pour une matrice carrée donnée, à chercher une matrice diagonale ou triangulaire qui lui est semblable.
	
	\section{Éléments propres}
	
	\subsection{Éléments propres d'un endomorphisme}
	
	Soit $u\in\mathcal{L}(E)$.
	
	\begin{dfn}
		\begin{itemize}
			\item On dit que $\lambda\in\K$ est \textbf{valeur propre} de $u$ s'il existe un vecteur \textit{non nul} $x\in E$ tel que $u(x)=\lambda x$.
			\item On dit que $x\in E$ est un \textbf{vecteur propre} de $u$ \textbf{associé à la valeur propre} $\lambda\in\K$ s'il est \textit{non nul} et vérifie $u(x)=\lambda x$.
			\item Si $\lambda\in\K$ est valeur propre de $u$, le \textbf{sous-espace propre associé à la valeur propre} $\lambda$ est : $$E_\lambda(u)=\ker(u-\lambda\Id_E)=\{x\in E\ : \ u(x)=\lambda x\}.$$
		\end{itemize}
	\end{dfn}
	
	\begin{rmq}
		\begin{itemize}
			\item Un vecteur propre n'est associé qu'à une seule valeur propre.
			\item Une droite vectorielle est stable par $u$ si, et seulement si, elle est engendrée par un vecteur propre de $u$. %IMPORTANT
			\item Si $\lambda\in\K$ est valeur propre de $u$, les vecteurs propres associés à la valeur propre $\lambda$ sont les vecteurs \textit{non nuls} de $E_\lambda(u)$. En particulier, si $E=\{0\}$, il n'y a aucun vecteur propre et donc aucune valeur propre.
			\item Plus généralement, on notera $E_\lambda(u)=\ker(u-\lambda\Id_E)$ pour tout $\lambda\in\K$. Il est non réduit à $\{0\}$ si, et seulement si, $\lambda$ est valeur propre de $u$.
			\item \textbf{Lien avec le noyau.} On a $E_0(u)=\ker(u)$, donc l'endomorphisme $u$ admet $0$ pour valeur propre si, et seulement s'il n'est pas injectif? Plus généralement, $\lambda$ est valeur propre de $u$ si, et seulement si, $u-\lambda\Id_E$ n'est pas injectif.
			\item \textbf{Lien avec l'image.} Tout sous-espace propre associé à une valeur propre non nulle est inclus dans $\text{Im}(u)$.\\ En effet, si $\lambda\ne 0$ et si $x\in E_\lambda(u)$, alors $x=\displaystyle u\left(\frac{x}{\lambda}\right)\in\text{Im}(u)$.
		\end{itemize}
	\end{rmq}
	
	\begin{exo}
		Pour $(a,b)\in\K\st\times\K$, quelles sont les valeurs propre de $au+b\Id_E$ ?
	\end{exo}
	
	\begin{met}
		Pour rechercher les éléments propres d'un endomorphisme $u$ (valeurs propres et sous-espaces propres associés), on peut résoudre l'équation $u(x)=\lambda x$, appelée équation aux éléments propres.
	\end{met}
	
	\paragraph{Endomorphismes remarquables}
	\begin{itemize}
		\item \textbf{Homothétie.} Tout vecteur non nul est vecteur propre de l'homothétie $\lambda\Id_E$ pour la valeur propre $\lambda$. Par suite, si $E\ne\{0\}$, cette homothétie admet $\lambda$ pour unique valeur propre et le sous-espace propre associé est $E$. Rappelons le résultat classique, reformulé ainsi : si tout vecteur non nul de $E$ est vecteur propre de $u$, alors $u$ est une homothétie.
		\item \textbf{Rotation.} Dans un plan euclidien, une rotation d'angle non multiple de $\pi$ n'a pas de valeur propre ni de vecteur propre.
		\item \textbf{Projection.} Soit $p\in\mathcal{L}(E)$ un projecteur. \\
		Si $\lambda$ est une valeur propre non nulle de $p$, alors $E_\lambda(p)\subset\text{Im}(p)$ donc $\lambda=1$ puisque $p$ coïncide avec l'identité sur son image. Il ne peut donc y avoir que deux valeurs propres, $0$ et $1$, avec : $$E_0(p)=\ker(p)\quad\text{et}\quad E_1(p)=\ker(p-\Id_E)=\text{Im}(p).$$
		Réciproquement, si $p\ne\Id_E$ et $p\ne 0$, on a $\ker(p)\ne\{0\}$ et $\text{Im}(p)\ne\{0\}$, donc $0$ et $1$ sont les valeurs propres de $p$.
		\item \textbf{Symétries.} Soit $F$ et $G$ supplémentaires dans $E$ et $s$ la symétrie par rapport à $F$ parallèlement à $G$. \\
		Rappelons que si $p$ est le projecteur sur $F$ parallèlement à $G$, on a $s=2p-\Id_E$. Ainsi, la relation $s(x)=\lambda x$ équivaut à $p(x)=\displaystyle\frac{\lambda+1}{2}x$. \\
		On déduit de ce qui a été prouvé sur les projections que si $s\ne\pm\Id_E$, alors $s$ admet pour valeur propres $1$ et $-1$ et que les sous-espaces propres associés sont respectivement $F$ et $G$.
		\item \textbf{Endomorphismes nilpotents} Un endomorphisme nilpotent $u$ n'admet que $0$ pour valeur propre. En effet, si $\lambda$ est une valeur propre de $u$ et $x$ un vecteur propre associé, on a $u(x)=\lambda x$ et pa une récurrence immédiate : $\forall k\in\N,\ u^k(x)=\lambda^kx$. \\
		En prenant $k$ tel que $u^k=0$,, cela donne $\lambda^k=0$ puisque $x\ne 0$ donc $\lambda=0$.\\
		Réciproquement, si $E\ne\{0\}$, alors $0$ est valeur propre de $u$ puisqu'un endomorphisme nilpotent est alors injectif. %retenir que $u$ nilpotent implique $Sp(u)\subset \{0\}$
	\end{itemize}
	
	\begin{ex}
		\begin{enumerate}
			\item Soit $E=\mathcal{C}^\infty(\R,\K)$, avec $\k=\R$ ou $\K=\C$. \\
			Tout $\lambda\in\K$ est valeur propre de l'endomorphisme $D$ de dérivation de $E$, et $E_\lambda(D)=\Vect(e_\lambda)$ où $e_\lambda:t\mapsto e^{\lambda t}$.
			\item Soit $E=\K^\N$ et $T\in\mathcal{L}(E)$ défini par :
			$$T((u_n)_{n\in\N})=(v_n)_{n\in\N}\quad\text{avec}\quad \forall n\in\N,\ v_n=u_{n+1}.$$
			\item Soit $E=\K[X]$ et $u\in\mathcal{L}(E)$ défini par $u(x)=XP$.
		\end{enumerate}
	\end{ex}
	
	\begin{exo}
		Soit $\varphi$ un isomorphisme de $E$ dans $F$ et $u$ un endomorphisme de $E$. Déterminer les éléments propres de $\varphi\circ u\circ\varphi\inv$ en fonction de ceux de $u$.
	\end{exo}
	
	\subsection{Propriétés}
	
	\begin{prop}%prop
		Si deux endomorphismes commutent, les sous-espaces propres de l'un sont stables par l'autre.
	\end{prop}
	
	Le résultat suivant est évident par définition d'un sous-espace propre.
	
	\begin{prop}%prop
		Si $F$ est un sous-espace vectoriel de $E$ stable par $u$, les valeurs propres de l'endomorphisme $u_F$ sont les valeurs propres de $E_\lambda(u)\cap F\ne\{0\}$. On a alors : $$E_\lambda(u_F)=E_\lambda(u)\cap F.$$
	\end{prop}
	
	\begin{prop}%prop
		Si $\lambda_1,\ldots,\lambda_p$ sont des valeurs propres de $u$ deux à deux distinctes, alors les sous-espaces propres associés $E_{\lambda_1}(u),\ldots,E_{\lambda_p}(u)$ sont en somme directe.
	\end{prop}
	\begin{proof}
		Appliquer le lemme de décomposition des noyaux à $P=(X-\lambda_1)\times\cdots\times(X-\lambda_p)$.
	\end{proof}
	
	\begin{exo}
		Redémontrer ce résultat :
		\begin{itemize}
			\item par récurrence sur $p$ ;
			\item en utilisant les polynômes d'interpolation de Lagrange.
		\end{itemize}
	\end{exo}
	
	\begin{cor}%cor
		Toute famille de vecteurs propres associés à deux valeurs propres deux à deux distinctes est libre.
	\end{cor}
	\begin{proof}
		Ce sont des vecteurs non nuls appartenant à des sous-espaces vectoriels en somme directe.
	\end{proof}
	
	\begin{ex}
		\begin{enumerate}
			\item S$\lambda_1,\ldots,\lambda_p$ sont des scalaires deux à deux distincts, alors la famille $(e_{\lambda_1},\ldots,e_{\lambda_p})$ est une famille libre de $\mathcal{C}^\infty(\R,\K)$ (où on a noté $e_{\lambda}:t\mapsto e^{\lambda t}$).
			\item Si $\lambda_1,\ldots,\lambda_p$ sont des scalaires deux à deux distincts, la famille de suites géométriques $((\lambda_1^n)_{n\in\N},\\ldots,(\lambda_p^n)_{n\in\N})$ est libre dans $\K^\N$. Ainsi, la famille $((\lambda^n)_{n\in\N})_{\lambda\in\K}$ est libre.
		\end{enumerate}
	\end{ex}
	
	\begin{exo}
		Pour tout $\alpha\in\R$, notons $f_\alpha$ la fonction définie sur $\R_+\st$ définie par $f_\alpha(t)=t^\alpha$. Montrer que la famille $(f_\alpha)_{\alpha\in\R}$ est libre.
	\end{exo}
	
	\begin{exo}
		Montrer que dans $\mathcal{C}^\infty(\R,\R)$, la famille constituée des fonctions $x\mapsto\cos(nx)$ (pour $n\in\N$) et $x\mapsto\sin(nx)$ (pour $n\in\N\st$) est libre. %penser à un opérateur de dérivation seconde
	\end{exo}
	
	\subsection{Cas de la dimension finie}
	
	\begin{cor}%cor
		Si $E$ est de dimension finie et si $\lambda_1,\ldots,\lambda_p$ sont des valeurs propres de $u$ deux à deux distinctes, alors : $$\sum_{i=1}^{p}\dim(E_{\lambda_i}(u))\leqslant\dim(E).$$
	\end{cor}
	
	\begin{dfn}
		Si $u$ est un endomorphisme d'un espace vectoriel de dimension finie, l'ensemble des valeurs propres de $u$ est le \textbf{spectre} de $u$, que l'on note $\Sp(u)$.
	\end{dfn}
	
	\begin{attention}
		Il arrive parfois que l'on parle du spectre d'un endomorphisme $u$ en dimension infinie : il s'agit alors de l'ensemble des scalaires $\lambda$ tels que $u-\lambda\Id_E$ soit non \textit{bijectif}, alors que les valeurs propres de $u$ sont les scalaires $\lambda$ tels que $u-\lambda\Id_E$ soit non injectif. Les éléments du spectre sont appelés les \textbf{valeurs spectrales} de $u$.
	\end{attention}
	
	\begin{ex}
		L'endomorphisme $P\mapsto XP$ de $\K[X]$ n'admet pas de valeur propre. En particulier, $0$ n'en est pas valeur propre, alors que $0$ en est valeur spectrale puisque cet endomorphisme est non bijectif (son image est incluse dans l'hyperplan d'équation $P(0)=0$).
	\end{ex}
	
	\begin{cor}%cor
		Le spectre d'un endomorphisme d'un espace vectoriel de dimension $n$ a au plus $n$ éléments distincts.
	\end{cor}
	
	\subsection{Éléments propres d'une matrice carrée}
	
	Soit $A\in\mathcal{M}_n(\K)$.
	
	\begin{dfn}
		Les \textbf{éléments propres} de $A$ sont les éléments propres de son endomorphisme canoniquement associé. Plus précisément :
		\begin{itemize}
			\item on dit que $\lambda\in\K$ est \textbf{valeur propre} de $A$ s'il existe $X\in\K^n$ \textit{non nul} tel que $AX=\lambda X$ ;
			\item on dit que $X\in\K^n$ est un \textbf{vecteur propre} de $A$ \textbf{associé à la valeur propre} $\lambda\in\K$ s'il est \textit{non nul} et vérifie $AX=\lambda X$ ;
			\item si $\lambda\in\K$ est valeur propre de $A$, le \textbf{sous-espace propre} de $A$ \textbf{associé à la valeur propre} $\lambda$ est :
			$$E_\lambda(A)=\ker(A-\lambda I_n)=\{X\in\K^n\ :\ AX=\lambda X\}\ ;$$
			\item l'ensemble des valeurs propres de $A$ est appelé le \textbf{spectre} de $A$ et noté $\Sp(A)$.
		\end{itemize}
	\end{dfn}
	
	\begin{rmq}
		Le scalaire $\lambda$ est valeur propre de $A$ si, et seulement si, $A-\lambda I_n$ est non inversible, c'est-à-dire si, et seulement si, $\rg(A-\lambda I_n)<n$.\\
		De plus, d'après le théorème du rang, $\dim(E_\lambda(A))=n-\rg(A-\lambda I_n)$.
	\end{rmq}
	
	\begin{exo}
		Soit $n\in\N$ tel que $n\geqslant 2$.
		\begin{enumerate}
			\item Déterminer les éléments propres de la matrice $J=\displaystyle\begin{pmatrix}
				1&\cdots&1\\
				\vdots&\ddots&\vdots\\
				1&\cdots&1
			\end{pmatrix}\in\mathcal{M}_n(\K)$.
			\item Soit $(\alpha,\beta)\in\K\times\K\st$. Déterminer les éléments propres de la matrice $A=\displaystyle\begin{pmatrix}
				\alpha&&(\beta)\\
				&\ddots&\\
				(\beta)&&\alpha
			\end{pmatrix}\in\mathcal{M}_n(\K)$.
		\end{enumerate}
	\end{exo}
	
	\begin{cor}%cor
		Une matrice $A\in\mathcal{M}_n(\K)$ admet au plus $n$ valeurs propres.
	\end{cor}
	
	\begin{met}
		Pour rechercher les éléments propres d'une matrice $A$, on peut résoudre l'équation $AX=\lambda X$, appelée équation aux éléments propres.
	\end{met}
	
	\begin{rmq}
		On verra comment, à l'aide du polynôme caractéristique ou d'un polynôme annulateur, obtenir les valeurs propres d'une matrice sans résolution de système.
	\end{rmq}
	
	\begin{ex}
		\begin{enumerate}
			\item Les valeurs propres d'une matrice triangulaire (supérieure ou inférieure) sont les éléments de sa diagonale.
			\item Soit $(a_0,\ldots,a_{n-1})\in\K^n$. Les valeurs propres de la transposée de la matrice compagnon du polynôme $P=X^n-a_{n-1}X^{n-1}-\cdots-a_1X-a_0$ : $$C_P^T=\begin{pmatrix}
				0&1&0&\cdots&0\\
				\vdots&\ddots&\ddots&\ddots&\vdots\\
				\vdots&\ddots&\ddots&\ddots&0\\
				0&\cdots&\cdots&0&1\\
				a_0&a_1&\cdots&a_{n-2}&a_{n-1}
			\end{pmatrix}\in\mathcal{M}_n(\K)$$ sont les racines de $P$ et les sous-espaces propres sont des droites.
		\end{enumerate}
	\end{ex}
	
	\begin{prop}%prop
		Soit $A$ une matrice représentant $u$ dans une base $\mathcal{B}=(e_1,\ldots,e_n)$.\\
		On a alors $\Sp(A)=\Sp(u)$ et, pour tout $\lambda\in\Sp(u)$ : $$x=\sum_{i=1}^{n}x_ie_i\in E_\lambda(u)\iff X=\begin{pmatrix}
			x_1\\
			\vdots\\
			x_n
		\end{pmatrix}\in E_\lambda(A).$$
	\end{prop}
	\begin{proof}
		Si $X$ est la matrice représentant un vecteur $x\in E$ dans la base $\mathcal{B}$, on a l'équivalence $u(x)=\lambda x\iff AX=\lambda X$.
	\end{proof}
	
	\begin{cor}{}%cor
		Deux matrices semblables ont le même spectre et les sous-espaces propres associés sont de même dimension.
	\end{cor}
	
	\begin{rmq}
		Plus précisément, si $B=PAP\inv$, alors pour tout $\lambda\in\Sp(B)$ : $$E_\lambda(B)=\{PX\ \mid\ X\in E_\lambda(A)\}.$$
	\end{rmq}
	
	\subsubsection*{Cas des matrices réelles considérées comme matrices complexes}
	
	Soit $A$ une matrice de taille $n\times n$ à coefficients réels. On peut considérer que $A$ appartient à $\mathcal{M}_n(\R)$ ou à $\mathcal{M}_n(\C)$.
	\begin{itemize}
		\item On obtient, dans le premier cas, les éléments propres réels, dans le second, les éléments propres complexes de $A$.
		\item On distinguera donc le spectre réel de $A$, noté $\Sp_\R(A)$, de son spectre complexe, noté $\Sp_\C(A)$.
	\end{itemize}
	
	\begin{prop}{}%prop
		Soit $A$ une matrice réelle et $\lambda\in\C$. On a :
		$$\lambda\in\Sp_\C(A)\iff\overline{\lambda}\in\Sp_C(A)\quad\text{et}\quad \forall X\in\C^n,\ X\in E_\lambda(A)\iff\overline{X}\in E_{\overline{\lambda}}(A).$$
	\end{prop}
	
	\begin{rmq}
		On peut démontrer de la même façon que si $A$ est une matrice complexe, alors $\Sp(\overline{A})=\overline{\Sp(A)}$ et que, pour tout $\lambda\in\Sp(A)$, les éléments de $E_\lambda(A)$ et de $E_{\overline{\lambda}}(\overline{A})$ sont conjugués.
	\end{rmq}
	
	\begin{rmq}
		Soit $A\in\mathcal{M}_n(\R)$.
		\begin{itemize}
			\item On a l'égalité $\Sp_\R(A)=\Sp_\C(A)\cap\R$. \\
			En effet, $\lambda$ est valeur propre de $A$ si, et seulement si, $A-\lambda I_n$ est non inversible, c'est-à-dire si, et seulement si, $\det(A-\lambda I_n)=0$. Or, pour $\lambda\in\R$, le déterminant de $A-\lambda I_n$ est le même que l'on considère cette matrice comme réelle ou complexe.
			\item En distinguant les deux sous-espaces propres réel et complexe de $A$ : $$E_\lambda^{(\R)}(A)=\{X\in\R^n\ :\ AX=\lambda X\}\quad\text{et}\quad E_\lambda^{(\C)}(A)=\{X\in\C^n\ :\ AX=\lambda X\},$$ on a $\dim_\R(E_\lambda^{(\R)}(A))=\dim_\C(E_\lambda^{(\C)}(A))$ puisque chacune de ces deux dimensions est égale à $n-\rg(A-\lambda I_n)$ d'après le théorème du rang.
		\end{itemize}
	\end{rmq}
	
	\begin{ex}
		Soit $A=\displaystyle\begin{pmatrix}
			0&-1\\
			1&0
		\end{pmatrix}\in\mathcal{M}_2(\R)$. On a $\Sp_\C(A)=\{i,-i\}$ et $\Sp_\R(A)=\emptyset$.
	\end{ex}
	
	\begin{rmq}
		Plus généralement, si $\L$ est un sous-corps de $\K$ et $A\in\mathcal{M}_n(\L)$, on a : $$\Sp_\L(A\subset)\Sp_\K(A)\quad\text{et, plus précisément,}\quad\Sp_\L(A)=\Sp_\K(A)\cap\L.$$
	\end{rmq}
	
	\subsection{Utilisation d'un polynôme annulateur}
	
	Soit $u\in\mathcal{L}(E)$ et $A\in\mathcal{M}_n(\K)$.
	
	\begin{prop}
		Soit $P$ un polynôme à coefficients dans $\K$ et $\lambda\in\K$.
		\begin{itemize}
			\item Si $x\in E_\lambda(u)$, on a $P(u)(x)=P(\lambda)x$.
			\item Si $X\in E_\lambda(A)$, on a $P(A)X=P(\lambda)X$.
		\end{itemize}
	\end{prop}
	\begin{proof}
		Récurrence immédiate pour la base canonique puis raisonnement par linéarité.
	\end{proof}
	
	[Ouvre le programme] "Y a que deux bibles : le programme officiel et le tout-en-un !"
	
	\begin{prop}%prop
		Si $P$ est un polynôme annulateur de $u$, alors toute valeur propre de $u$ est racine de $P$. On a le même résultat pour une matrice $A\in\mathcal{M}_n(\K)$.
	\end{prop}%prop précédente et pseudo intégrité
	
	\begin{ex}
		\begin{enumerate}
			\item On retrouve ainsi que le spectre d'un projecteur $p$ est inclus dans $\{0,1\}$ puisque $p$ est annulé par $X^2-X$. De même, celui d'une symétrie est inclus dans $\{-1,1\}$ à l'aide du polynôme $X^2-1$.
			\item $A=\displaystyle\begin{pmatrix}
				0&0&1\\
				1&0&0\\
				0&1&0
			\end{pmatrix}.$
		\end{enumerate}
	\end{ex}
	
	\begin{attention}
		Lorsque l'on dispose d'un polynôme annulateur $P$ de $u\in\mathcal{L}(E)$, toutes les valeurs propres de $u$ sont racines de $P$, mais certaines racines de $P$ peuvent ne pas être valeur propre de $u$.\\
		Ainsi, le polynôme $X^2-X=X(X-1)$ est un polynôme annulateur de tout projecteur, dont $\Id_E$, alors que $0$ n'est pas valeur propre de l'identité.
	\end{attention}
	
	\begin{notation}
		Lorsqu'un endomorphisme $u$ admet un polynôme annulateur non nul, son polynôme minimal est noté $\pi_u$ ou $\mu_u$. De même, le polynôme minimale d'une matrice carrée $A$ est noté $\pi_A$ ou $\mu_A$.
	\end{notation}
	
	\begin{prop}%prop
		Si un endomorphisme $u$ admet un polynôme annulateur non nul, ses valeurs propres sont exactement les racines de $\pi_u$.
	\end{prop}
	\begin{proof}
		Écrire $\pi_u=(X-\lambda)P$, et prendre $x$ tel que $P(u)(x)\ne 0$.
	\end{proof}
	
	\begin{ex}
		Polynôme minimal et valeurs propres d'une matrice compagnon : $$C_P=\begin{pmatrix}
			0&\cdots&\cdots&0&a_0\\
			1&\ddots&\ddots&\vdots&\vdots\\
			0&\ddots&\ddots&\vdots&\vdots\\
			\vdots&\ddots&\ddots&0&a_{n-2}\\
			0&\cdots&0&1&a_{n-1}
		\end{pmatrix}.$$
	\end{ex}
	
	\section{Polynôme caractéristique}
	
	\subsection{Polynôme caractéristique d'une matrice}
	
	Soit $A\in\mathcal{M}_n(\K)$.\\
	Par définition, $\lambda\in\K$ est valeur propre de $A\in\mathcal{M}_n(\K)$ so, et seulement si, la matrice $A-\lambda I_n$ est non inversible, donc si, et seulement si, $\det(\lambda I_n-A)=0$.\\
	Comme le déterminant d'une matrice est une fonction polynomiale en ses coefficients et que les coefficients de $\lambda I_n-A$, on en déduit que $\det(\lambda I_n-A)$ est polynomiale en $\lambda$. Cela justifie la définition suivante.
	
	\begin{dfn}
		On appelle \textbf{polynôme caractéristique} de $A$, et l'on note $\chi_A$, l'unique polynôme tel que : $$\forall\lambda\in\K,\ \chi_A(\lambda)=\det(\lambda I_n-A).$$
	\end{dfn}
	
	\begin{rmq}
		\begin{itemize}
			\item Il y a bien unicité d'un tel polynôme. En effet, $\K$ étant un sous-corps de $\C$, il contient $\Q$, qui est infini, et donc deux polynôme ayant la même fonction polynomiale associée sont égaux.
			\item Par abus, on note parfois $\chi_A(X)=\det(XI_n-A)$ que l'on peut considérer comme le déterminant d'une matrice à coefficients dans le corps $\K(X)$ des fractions rationnelles. il s'agit d'un léger abus dans le cadre du programme. En effet, on a défini uniquement le déterminant d'une matrice à coefficients dans un sous-corps de $\C$.
			\item Si $A\in\mathcal{M}_n(\R)$, son polynôme caractéristique $\chi_A$ est identique que l'on considère $A$ comme une matrice réelle ou complexe.
		\end{itemize}
	\end{rmq}
	
	\begin{ex}
		On a $\chi_A=\chi_{A^T}$.
	\end{ex}
	
	\begin{prop}%prop
		Le polynôme caractéristique $\chi_A$ est un polynôme unitaire de degré $n$ et l'on a : $$\chi_A=X^n-\Tr(A)X^{n-1}+\cdots+(-1)^n\det(A).$$
	\end{prop}
	\begin{proof}
		Utiliser la formule des signatures.
	\end{proof}
	
	\begin{exo}
		Soit $A\in\mathcal{A}\in\mathcal{GL}_n(\K)$. Pour $\lambda\ne0$, donner une relation entre $\chi_A(\lambda)$ et $\chi_{A\inv}(\lambda\inv)$. En déduire une expression du coefficient de $X$ dans $\chi_A$ en fonction de la trace de la comatrice de $A$.
	\end{exo}
	
	\begin{met}
		Le polynôme caractéristique de $A\in\mathcal{M}_2(\K)$ est : $$X^2-\Tr(A)X+\det(A)$$ On utilise impérativement cette expression quand on travaille dans $\mathcal{M}_2(\K)$.
	\end{met}
	
	\begin{prop}%prop
		Deux matrices semblables ont le même polynôme caractéristique.
	\end{prop}
	
	\begin{ex}
		Pour tout $(A,B)\in\mathcal{M}_n(\K)^2$, on a $\chi_{AB}=\chi_{BA}$.
	\end{ex}
	
	\subsection{Exemples de calcul de polynômes caractéristiques}
	
	\begin{ex}
		\begin{itemize}
			\item Si $M$ est une matrice triangulaire supérieure par blocs de la forme $\displaystyle\begin{pmatrix}
				A&C\\
				0&D
			\end{pmatrix}$, alors son polynôme caractéristique est $\chi_M=\chi_A\chi_D$. Ce résultat se généralise à toute matrice triangulaire (supérieure ou inférieure) par blocs : le polynôme caractéristique est le produit des polynômes caractéristiques des blocs diagonaux.
			\item $A=\displaystyle\begin{pmatrix}
				1&0&0\\
				0&\cos(\theta)&-\sin(\theta)\\
				0&\sin(\theta)&\cos(\theta)
			\end{pmatrix}$.
		\end{itemize}
	\end{ex}
	
	\begin{exo}
		Calculer le polynôme caractéristique de la matrice par blocs $B=\displaystyle\begin{pmatrix}
			0&A\\
			A&0
		\end{pmatrix}$ et fonction de celui de $A\in\mathcal{M}_n(\K)$.
	\end{exo}
	
	\begin{prop}%prop
		Si $A\in\mathcal{M}_n(\K)$ est triangulaire de diagonale $(\alpha_1,\ldots,\alpha_n)$, alors son polynôme caractéristique est : $$\chi_A=\prod_{i=1}^{n}(X-\alpha_i).$$
	\end{prop}
	
	\begin{ex}
		Soit $(a_0,\ldots,a_{n-1})\in\K^n$. Polynôme caractéristique de la matrice compagnon :
		$$C_P=\begin{pmatrix}
			0&\cdots&\cdots&0&a_0\\
			1&\ddots&\ddots&\vdots&\vdots\\
			0&\ddots&\ddots&\vdots&\vdots\\
			\vdots&\ddots&\ddots&0&a_{n-2}\\
			0&\cdots&0&1&a_{n-1}
		\end{pmatrix}.$$
	\end{ex}
	
	\begin{met}
		Comme on l'a fait dans l'exemple précédent, il parfois est plus intéressant de calculer $\det(A-\lambda I_n)=(-1)^n\chi_{_A}(\lambda)$ plutôt que $\det(\lambda I_n-A)=\chi_{_A}(\lambda)$.
	\end{met}
	
	\subsection{Polynôme caractéristique et spectre}
	
	\begin{prop}%prop
		Un scalaire $\lambda\in\K$ est une valeur propre de $A$ si, et seulement s'il est racine du polynôme caractéristique de $A$.
	\end{prop}
	
	\begin{rmq}
		Comme on sait que le polynôme caractéristique d'une matrice de taille $n$ est de degré $n$, cela permet donc de retrouver qu'une matrice carrée de taille $n$ a au plus $n$ valeurs propres.
	\end{rmq}
	
	\begin{ex}
		Les valeurs propres d'une matrice triangulaire sont ses coefficients diagonaux.
	\end{ex}
	
	\begin{met}
		Puisque les racines du polynôme caractéristique sont les valeurs propres, il est important d'obtenir le polynôme caractéristique sous forme factorisée. On réalise dans les cas concret cet objectif en calculant le déterminant de $A-\lambda I_n$ par opérations élémentaires afin (suivant les cas) :
		\begin{itemize}
			\item de faire apparaître des facteurs communs dans les lignes ou les colonnes ;
			\item de se ramener à une matrice triangulaire, au moins en partie.
		\end{itemize} On pourra notamment utiliser des combinaisons linéaires qui traduisant des relations de liaison entre les colonnes, ou encore passer à la transposée si ces relations de liaison sont sur les lignes.
	\end{met}
	
	\begin{ex}
		Déterminons le spectre de $A=\begin{pmatrix}
			3&5&-6\\
			4&7&-9\\
			3&6&-7
		\end{pmatrix}$ en remarquant que la somme des colonnes...
	\end{ex}
	
	\begin{rmq}
		Soit $A\in\mathcal{M}_n(\K)$, avec $n\in\N\st$.\\
		En utilisant la proposition précédente, on voit que :
		\begin{itemize}
			\item si $\K=\C$, alors $A$ a au moins une valeur propre, d'après le théorème de d'Alembert-Gauss ;
			\item si $\K=\R$ et si $n$ est impair, alors $A$ a au moins une valeur propre d'après le théorème des valeurs intermédiaires.
		\end{itemize}
	\end{rmq}
	
	\subsection{Polynôme caractéristique d'un endomorphisme}
	
	Dans tout cette section, on suppose $E$ de dimension finie $n$ et l'on considère un endomorphisme $u$ de $E$.
	
	\begin{dfn}
		On appelle \textbf{polynôme caractéristique} de $u$, et l'on note $\chi_u$, l'unique polynôme vérifiant : $$\forall \lambda\in\K,\ \chi_u(\lambda)=\det(\lambda\Id_E-u),$$c'est-à-dire le polynôme caractéristique de sa matrice dans n'importe quelle base. 
	\end{dfn}
	
	En particulier, si $u$ est l'endomorphisme canoniquement associé à la matrice $A$, alors $\chi_u=\chi_A$.
	
	\begin{attention}
		La notion de polynôme caractéristique d'un endomorphisme n'a aucun sens si $E$ n'est pas de dimension finie.
	\end{attention}
	
	\begin{ex}
		Soit $u\in\mathcal{L}(E)$ de rang $1$. Son polynôme caractéristique est : $$\chi_u(X)=X^{n-1}(X-\Tr(u)).$$
	\end{ex}
	
	\begin{prop}{}%prop
		Un scalaire $\lambda\in\K$ est une valeur propre de $u$ si, et seulement si, c'est une racine du polynôme caractéristique de $u$.
	\end{prop}
	
	\begin{rmq}
		\begin{itemize}
			\item On retrouve le fait qu'un endomorphisme d'un espace vectoriel de dimension $n$ a au plus $n$ valeurs propres distinctes.
			\item Soit $E$ un $\K$-espace vectoriel de dimension finie non nulle.
			\begin{itemize}[label=$\star$]
				\item Si $\K=\C$, alors $u$ a au moins une valeur propre.
				\item Si $\K=\R$ et si $\dim(E)$ est impaire, alors $u$ a au moins une valeur propre.
			\end{itemize}
		\end{itemize}
	\end{rmq}
	
	\begin{attention}
		En dimension infinie, un endomorphisme d'un $\C$-espace vectoriel n'admet pas nécessairement de valeur propre.
	\end{attention}
	
	\begin{prop}{}%prop
		Si $F$ est un sous-espace vectoriel de $E$ stable par $u$, alors le polynôme caractéristique de $u_F$ divise le polynôme caractéristique de $u$.
	\end{prop}
	
	\begin{rmq}
		\begin{itemize}
			\item On retrouve l'inclusion $\Sp(u_F)\subset\Sp(u)$.
			\item Le polynôme caractéristique d'un endomorphisme $u$ stabilisant les sous-espaces vectoriels d'une décomposition $E=E_1\oplus\cdots\oplus E_p$ est : $$\chi_u=\prod_{i=1}^{p}\chi_{u_i},$$ où, pour tout $i\in\llbracket1,p\rrbracket$, $u_i$ est l'endomorphisme induit par $u$ sur $E_i$.
		\end{itemize}
	\end{rmq}
	
	\subsection{Ordre de multiplicité d'une valeur propre}
	
	Soit $u\in\mathcal{L}(E)$, où $E$ est toujours supposé de dimension finie.
	
	\begin{dfn}
		On appelle \textbf{ordre de multiplicité} d'une valeur propre $\lambda$ de $u$, son ordre de multiplicité en tant que racine du polynôme caractéristique de $u$.\\
		On définit de même l'ordre de multiplicité des valeurs propres d'une matrice carrée.
	\end{dfn}
	
	\begin{rmq}
		En particulier, une valeur propre de $u$ est dite simple, double, triple... si c'est une racine simple, double, triple... du polynôme caractéristique de $u$.
	\end{rmq}
	
	\begin{notation}
		Dans la suite, l'ordre de multiplicité de $\lambda$ sera noté $m(\lambda)$.
	\end{notation}
	
	\begin{attention}
		La notion d'ordre de multiplicité n'a pas de sens en dimension infinie.
	\end{attention}
	
	\begin{prop}{}%prop
		Pour tout $\lambda\in\Sp(u)$, on a : $$1\leqslant\dim(E_\lambda(u))\leqslant m(\lambda).$$
	\end{prop}
	\begin{proof}
		$(X-\lambda)^{\dim(E_\lambda(u))}$ divise $\chi_u$ d'après la proposition précédente.
	\end{proof}
	
	\begin{ex}
		\begin{enumerate}
			\item Si $u$ est de rang inférieur ou égal à $r$, le polynôme $\chi_u$ est divisible par $X^{n-r}$, puisque le noyau $\ker(u)=E_0(u)$ est de dimension supérieure ou égale à $n-r$.
			\item Ainsi, le polynôme caractéristique de la matrice : 
			$$A=\begin{pmatrix}
				&&&\alpha_1\\
				&(0)&&\vdots\\
				&&&\alpha_{n-1}\\
				\alpha_1&\cdots&\alpha_{n-1}&0
			\end{pmatrix}$$ de range inférieur ou égal à $2$ est de la forme : $$\chi_A=X^{n-2}(X^2+aX+b).$$
		\end{enumerate}
	\end{ex}
	
	\begin{cor}{}%cor
		Si $\lambda$ est valeur propre simple de $u$, alors $\dim(E_\lambda(u))=1$.
	\end{cor}
	
	\begin{rmq}
		Le spectre de $u$ est, par définition, l'ensemble des racines de $\chi_u$ dans $\K$. Comme dans le cas des polynômes, on distinguera soigneusement les notions d'ensemble er de liste des valeurs propres de $u$.
		\begin{itemize}
			\item L'ensemble des valeurs propres est le spectre de $u$. S'il est égal à $\{\lambda_1,\lambda_p\}$, avec $\lambda_1,\ldots,\lambda_p$ deux à deux distincts, alors on a : $$\chi_u=Q\prod_{i=1}^{p}(X-\lambda_i)^{m(\lambda_i)}$$ où $Q\in\K[X]$ n'a pas de racine dans $\K$ et les $m(\lambda_i)$ sont les ordres de multiplicité.
			\item Une liste des valeurs propres comptées avec ordre de multiplicité est une famille de scalaires répétant les valeurs propres avec leurs multiplicités. Une telle liste est unique à l'ordre près et si $(\mu_1,\ldots,\mu_s)$ en est une, alors on a : $$\chi_u=Q\prod_{i=1}^{s}(X-\mu_i)$$ où $Q\in\K[X]$ n'a pas de racine dans $\K$.
			\item Bien sûr, la liste $(\mu_1,\ldots,\mu_s)$ est formée des $\lambda_1,\ldots,\lambda_p$ répétés autant de fois que leur multiplicité, c'est-à-dire, à l'ordre près :
			$$(\mu_1,\ldots,\mu_s)=(\underbrace{\lambda_1,\ldots,\lambda_1}_{m(\lambda_1)\text{ fois}},\ldots\ldots,\underbrace{\lambda_p,\ldots,\lambda_p}_{m(\lambda_p)\text{ fois}}).$$
			On a donc : $$s=m(\lambda_1)+\cdots+m(\lambda_p).$$
			\item Si $\chi_u$ est scindé, alors on a $s=n$ et : $$\chi_u=\prod_{k=1}^{p}(X-\lambda_k)^{m(\lambda_k)}=\prod_{i=1}^{n}(X-\mu_i).$$
		\end{itemize}
	\end{rmq}
	
	\begin{ex}
		La diagonale d'une matrice triangulaire constitue une liste de ses valeurs propres, comptées avec leur ordre de multiplicité.
	\end{ex}
	
	\begin{prop}{}%prop
		Si $\chi_u=\displaystyle\prod_{i=1}^{n}(X-\mu_i)$ est scindé, alors : $$\Tr(u)=\sum_{i=1}^{n}\mu_i\quad\text{et}\quad\det(u)=\prod_{i=1}^{n}\mu_i.$$
	\end{prop}
	\begin{proof}
		Relations coefficients-racines.
	\end{proof}
	
	\paragraph{Autre formulation} Si $\Sp(u)=\{\lambda_1,\lambda_p\}$ avec $\lambda_1,\ldots,\lambda_p$ deux à deux distincts : $$\Tr(u)=\sum_{i=1}^{p}m(\lambda_i)\lambda_i\quad\text{et}\quad\det(u)=\prod_{i=1}^{p}\lambda_i^{m(\lambda_i)}.$$
	
	\subsection{Théorème de Cayley-Hamilton}
	
	On a vu l'intérêt d'obtenir des polynômes annulateurs dans la recherche du spectre d'un endomorphisme et l'on prouvé leur existence en dimension finie.\\
	Le théorème de Cayley-Hamilton ci-dessous affirme que le polynôme caractéristique est une polynôme annulateur, ce qui nous permettra d'obtenir en pratique un polynôme annulateur.\\
	La démonstration de ce théorème n'est pas exigible.
	
	\begin{exo}
		Soit $x$ un vecteur non nul de $E$.
		\begin{enumerate}
			\item Montrer qu'il existe un plus petit sous-espace vectoriel de $E$, noté $E_x$, stable par $u$ et contenant $x$.
			\item Prouver que si $E_x$ est de dimension $p$, alors $\mathcal{B}_x=(x,u(x),\ldots,u^{p-1}(x))$ est une base de $E_x$.
			\item Montrer que la matrice dans $\mathcal{B}_x$ de l'endomorphisme $u_x$ induit par $u$ sur $E_x$ est une matrice compagnon.
		\end{enumerate}
	\end{exo}
	
	\begin{thm}[Théorème de Cayley-Hamilton]%thm
		Le polynôme caractéristique de $u\in\mathcal{L}(E)$ annule $u$, c'est-à-dire $\chi_u(u)=0$.\\
		De même, si $A\in\mathcal{M}_n(\K)$, alors $\chi_A(A)=0$.
	\end{thm}
	\begin{proof}
		Des exemples précédents montrent que $\pi_A=\chi_A$ pour une matrice compagnon. L'exercice précédent donner alors $\chi_u(u)(x)=0$ pour tout $x\in E$.
	\end{proof}
	
	\begin{attention}
		Comme il fait intervenir le polynôme caractéristique de $u\in\mathcal{L}(E)$, le théorème de Cayley-Hamilton n'a de sens que si $E$ est de dimension \textbf{finie}.
	\end{attention}
	
	\begin{exo}
		On reprend la matrice $A=\begin{pmatrix}
			3&5&-6\\
			4&7&-9\\
			3&6&-7
		\end{pmatrix}$. Montrer que $A$ est inversible et calculer $A\inv$ en fonction de $I_3$, $A$ et $A^2$.
	\end{exo}
	
	\begin{cor}{}%cor
		Le polynôme minimal d'un endomorphisme d'un espace vectoriel de dimension finie $n$ ou d'une matrice de $\mathcal{M}_n(\K)$ est de degré inférieur ou égal à $n$.
	\end{cor}
	
	\begin{exo}
		Que penser de ces deux démonstrations de la version matricielle du théorème de Cayley-Hamilton ?
		\begin{enumerate}
			\item Comme $\chi_A(X)=\det(XI_n-A)$, on a $\chi_A(A)=\det(AI_n-A)=\det(0)=0$.
			\item On utilise la relation : $$\forall N\in\mathcal{M}_n(\K),\ N\text{Com}^T(N)=\det(N)I_n,$$ que l'on applique à $N=XI_n-A$ à coefficients dans le corps $\K(X)$. La matrice $\text{Com}^T(N)$ est à coefficients polynomiaux : on peut donc aussi la voir comme un polynôme de degré $n-1$ à coefficients dans $\mathcal{M}_n(\K)$ soit : $$\text{Com}^T(N)=X^{n-1}A_{n-1}+X^{n-2}A_{n-2}+\cdots+XA_1+A_0.$$ Si $\det(N)=\chi_A(X)=a_nX^n+a_{n-1}X^{n-1}+\cdots+a_1X+a_0$, on identifie les deux polynômes, on élimine les $A_i$ et... on trouve : $$a_nA^n+a_{n-1}A^{n-1}\cdots+a_1A_1+a_0I_n=0.$$
		\end{enumerate}
	\end{exo}
	
	\section{Diagonalisation}
	
	Dans cette section, $u$ est un endomorphisme de $E$, où $E$ est supposé de dimension finie $n$, et $A\in\mathcal{M}_n(\K)$.
	
	\subsection{Définitions et premières propriétés}
	
	\begin{dfn}
		\begin{itemize}
			\item L'endomorphisme $u$ est dit \textbf{diagonalisable} s'il existe une base de $E$ dans laquelle la matrice de $u$ est diagonale.\\
			Une telle base s'appelle une \textbf{base de diagonalisation} de $u$.
			\item La matrice $A$ est dite \textbf{diagonalisable} si elle est semblable à une matrice diagonale, c'est-à-dire s'il existe $D$ diagonale et $P$ inversible telles que $A=PDP\inv$.
		\end{itemize}
	\end{dfn}
	
	\paragraph{Reformulation} Autrement dit, $u$ est diagonalisable si, et seulement s'il existe une base de $E$ constituée de vecteurs propres de $u$.
	
	\begin{prop}%prop
		Si $A$ représente $u$, elle est diagonalisable si, et seulement si, $u$ est diagonalisable.
	\end{prop}
	
	\begin{cor}%cor
		La matrice $A$ est diagonalisable si, et seulement si, son endomorphisme canoniquement associé est diagonalisable.
	\end{cor}
	
	\begin{terminologie}
		\begin{itemize}
		\item \textbf{Diagonaliser} un endomorphisme de $E$ signifie trouver une base de $E$ dans laquelle la matrice de $u$ soit diagonale, c'est-à-dire une base constituée de vecteurs propres de $u$.
		\item \textbf{Diagonaliser} une matrice signifie diagonaliser son endomorphisme canoniquement associé, ce qui revient à trouver une matrice $D$ diagonale et une matrice $P$ inversible telles que $A=PDP\inv$.
	\end{itemize}
	\end{terminologie}
	
	\begin{met}
		Soit $A\in\mathcal{M}_n(\K)$ diagonalisable et $(E_1,\ldots,E_n)$ une base de $\K^n$ constituée de vecteurs propres de $A$ associés aux valeurs propres $\lambda_1,\ldots,\lambda_n$. Si on pose $p$ la matrices dont les colonnes sont $E_1,\ldots,E_n$, alors : $$A=PDP\inv \text{  ou  }D=\text{Diag}(\lambda_1,\ldots,\lambda_n)$$
	\end{met}
	
	Exercice Moulin.Red.23
	
	\begin{ex}
		Une matrice est diagonalisable si, et seulement si, sa transposée est diagonalisable.
	\end{ex}
	
	\begin{prop}%prop
		Si $E$ est de dimension $n$ et si $u\in\mathcal{L}(E)$ possède $n$ valeurs propres distinctes, alors $u$ est diagonalisable et chaque sous-espace propre est de dimension $1$.\\
		De même, si $A\in\mathcal{M}_n(\K)$ possède $n$ valeurs propres distinctes, alors $A$ est diagonalisable et chaque sous-espace propre est de dimension $1$.
	\end{prop}
	
	\begin{ex}
		La matrice $\begin{pmatrix}
			1&2&3&\\
			0&4&5\\
			0&0&6
		\end{pmatrix}$ est diagonalisable.
	\end{ex}
	
	\begin{attention}
		Il ne s'agit bien que d'une condition \textit{suffisante} pour que $u$ soit diagonalisable. La proposition suivante donne une caractérisation, qui peut être utile dans le cas où le nombre de valeurs propres est strictement inférieur à la dimension de $E$.
	\end{attention}
	
	\begin{prop}%prop
		Les propriétés suivantes sont équivalentes :
		\begin{enumerate}[label=(\roman*)]
			\item l'endomorphisme $u$ est diagonalisable ;
			\item $\displaystyle\bigoplus_{\lambda\in\Sp(u)}E_\lambda(u)=E$ ;
			\item $\displaystyle\sum_{\lambda\in\Sp(u)}\dim(E_\lambda(u))=n$.
		\end{enumerate}
	\end{prop}
	
	\begin{cor}%cor
		Il y a équivalence entre :
		\begin{enumerate}[label=(\roman*)]
			\item la matrice $A$ est diagonalisable ;
			\item $\displaystyle\bigoplus_{\lambda\in\Sp(A)}E_\lambda(A)=\K^n$ ;
			\item $\displaystyle\sum_{\lambda\in\Sp(A)}\dim(E_\lambda(A))=n$.
		\end{enumerate}
	\end{cor}
	
	\begin{exo}
		La matrice $A=\begin{pmatrix}
			0&3&2\\
			2&5&2\\
			2&3&0
		\end{pmatrix}$ est-elle diagonalisable ? Si c'est le cas, fournir une matrice $P\in\mathcal{GL}_3(\K)$ telle que $P\inv AP$ soit diagonale. 
	\end{exo}
	
	\begin{ex}
		\begin{enumerate}
			\item Un projecteur $p$ de $E$ est diagonalisable.
			\begin{itemize}
				\item Si $p=0$ ou $p=\Id_E$, c'est immédiat.
				\item Sinon, on sait que $\Sp(p)=\{0,1\}$ et l'on a $E=E_0(p)\oplus E_1(p)$.
			\end{itemize}
			Un exemple suivant montrera comment éviter de traiter ces cas particuliers.
			\item Si $u\in\mathcal{L}(E)$ admet $\lambda$ pour unique valeur propre, alors $u$ est diagonalisable si, et seulement si, le sous-espace propre associé est égal à $E$, c'est-à-dire si, et seulement si, $u=\lambda\Id_E$.\\
			De même, si $A\in\mathcal{M}_n(\K)$ admet $\lambda$ pour unique valeur propre $\lambda$, alors elle est diagonalisable si, et seulement si, son endomorphisme canoniquement associé est égal à $\lambda\Id_{\K^n}$, c'est-à-dire si, et seulement si, $A=\lambda I_n$.
			\item En particulier, un endomorphisme nilpotent (respectivement une matrice nilpotente) n'est diagonalisable que s'il (elle) est nul(le). %orthogonalité des notions de nilpotence et de diagonalisabilité
			\item De même, si $A\in\mathcal{M}_n(\K)$ est triangulaire de coefficients diagonaux tous égaux à $\lambda$, son polynôme caractéristique est égal à $(X-\lambda)^n$. La matrice $A$ admet donc pour unique valeur propre $\lambda$, et n'est donc diagonalisable que si $A=\lambda I_n$.
		\end{enumerate}
	\end{ex}
	
	\begin{rmq}
		Les sou-espaces propres étant en somme directe, la somme de leurs dimensions est majorée par la dimension de $E$.
	\end{rmq}
	
	\begin{met}
		Pour montrer qu'un endomorphisme d'un espace vectoriel de dimension $n$, ou une matrice de taille $n$, est diagonalisable, il suffit de montrer que la somme des dimensions de ses sous-espaces propres est supérieure ou égale à $n$.
	\end{met}
	
	\begin{ex}
		Cherchons une condition nécessaire et suffisante pour que la matrice $$A(k)=\begin{pmatrix}
			0&1&0&0\\
			1&k&1&1\\
			0&1&0&0\\
			0&1&0&0
		\end{pmatrix}$$ soit diagonalisable.
	\end{ex}
	
	\begin{rmq}
		Si $E$ est somme de sous-espaces vectoriels stables $E_1,\ldots,E_p$ sur chacun desquels $u$ induit une homothétie, alors $u$ est diagonalisable.\\
		En effet, en prenant dans chaque $E_i$ des vecteurs en constituant une base, on obtient une famille génératrice de $E$ constituée de vecteurs propres, de laquelle on peut extraire une base.
	\end{rmq}
	
	\begin{ex}
		\begin{enumerate}
			\item Si $p$ est un projecteur de $E$, on sait que l'on a $E=E_0(p)\oplus E_1(p)$, avec $p_{E_0(p)}=0$ et $p_{E_1(p)}=\Id_{E_1(p)}$, donc $p$ est diagonalisable.
			\item De même, toute symétrie $s$ de $E$ est diagonalisable puisque $E=E_{-1}(s)\oplus E_1(s)$.
			\item L'endomorphisme $f\in\mathcal{L}(\mathcal{M}_n(\K))$ défini par $f(A)=\Tr(A)I_n+A$ est diagonalisable.
		\end{enumerate}
	\end{ex}
	
	\subsection{Utilisation du polynôme caractéristique}
	
	Les valeurs propres étant les racines du polynôme caractéristique, on peut reformuler une proposition précédente.
	
	\begin{prop}%prop
		Si le polynôme caractéristique de $u$ (respectivement de $A$) est scindé à racines simples, alors $u$ (respectivement $A$) est diagonalisable.
	\end{prop}
	
	\begin{terminologie}
		Un polynôme scindé à racines simples est dit aussi \textbf{simplement scindé}.
	\end{terminologie}
	
	\begin{attention}
		Il faut bien noter que le résultat précédent ne fournit qu'une condition \textit{suffisante} pour que $u$ soit diagonalisable comme le montre, par exemple, le cas d'une homothétie.\\	Le théorème qui suit donne une condition nécessaire et suffisante.
	\end{attention}
	
	\begin{thm}%thm
		Pour que $u$ soit diagonalisable, il faut et il suffit qu'il vérifie les deux conditions suivantes :
		\begin{itemize}
			\item son polynôme caractéristique est scindé sur $\K$ ;
			\item pour toute valeur propre de $u$, la dimension du sous-espace propre associé est égale à l'ordre de multiplicité de cette valeur propre, c'est-à-dire :
			$$\forall\lambda\in\Sp(u),\ \dim(E_\lambda(u))=m(\lambda).$$
		\end{itemize}
		On a le même résultat pour une matrice $A\in\mathcal{M}_n(\K)$.
	\end{thm}
	\begin{proof}
		On écrit $\chi_u=Q\displaystyle\prod_{\lambda\in\Sp(u)}(X-\lambda)^{m(\lambda)}$, où $Q$ n'a pas de racine dans $\K$, et on utilise le fait que $\dim(E_\lambda(u))\leqslant m(\lambda)$, pour tout $\lambda\in\Sp(u)$.
	\end{proof}
	
	\begin{rmq}
		Si $\lambda$ est valeur propre simple, on a toujours $\dim(E_\lambda(u))=m(\lambda)$ car : $$E_\lambda(u)\ne\{0\}\quad\text{et}\quad\dim(E_\lambda(u))\leqslant m(\lambda)=1.$$
		Dans la caractérisation précédente, on peut donc se limiter aux sous-espaces propres associés à des valeurs propres multiples.
	\end{rmq}
	
	\begin{met}
		Pour savoir si un endomorphisme $u$ est diagonalisable, on peut calculer son polynôme caractéristique.
		\begin{itemize}
			\item S'il n'est pas scindé, alors $u$ n'est pas diagonalisable.
			\item Sinon, pour toute valeur propre multiple, on compare $\dim(E_\lambda(u))$ et $m(\lambda)$, par exemple en calculant le rang de $u-\lambda\text{Id}_E$.
		\end{itemize}
	\end{met}
	
	\begin{ex}
		\begin{enumerate}
			\item Soit $A=\begin{pmatrix}
				1&a&b\\
				0&1&c\\
				0&0&-1
			\end{pmatrix}_in\mathcal{M}_3(\K)$. Donner une condition nécessaire et suffisante portant sur $(a,b,c)\in\K^3$ pour que la matrice $A$ soit diagonalisable.
			\item L'endomorphisme $v:P\mapsto (1-X^2)P'+nXP$ de $\R[X]$ induit un endomorphisme diagonalisable de $\R_n[X]$.
		\end{enumerate}
	\end{ex}
	
	\subsection{Utilisation d'un polynôme annulateur}
	
	A lui seul, le polynôme caractéristique ne peut permettre de montrer que $u$ est diagonalisable que s'il est scindé à racines simples. Sinon, on a besoin d'étudier en plus la dimension de ses sous-espaces propres.\\
	Le théorème suivant donne une condition nécessaire et suffisante pour être diagonalisable, en utilisant cette fois-ci les polynômes annulateurs.
	
	\begin{thm}%thm
		L'endomorphisme $u$ est diagonalisable si, et seulement s'il possède un polynôme annulateur scindé à racines simples.\\
		On a le même résultat si $A\in\mathcal{M}_n(\K)$.
	\end{thm}
	\begin{proof}
		Considérer le polynôme $P=\displaystyle\prod_{\lambda\in\Sp(u)}(X-\lambda)$ dans le sens direct, et utiliser le lemme de décomposition des noyaux pour le sens réciproque.
	\end{proof}
	
	\begin{ex}
		\begin{enumerate}
			\item On retrouve que les projecteurs et les symétries sont diagonalisables puisqu'ils sont annulés respectivement par les polynômes $X^2-X$ et $X^1-1$ scindés à racines simples.
			\item Une matrice carré est diagonalisable si, et seulement si, sa transposée l'est.
			\item Soit $\sigma\in\mathcal{S}_n$. La matrice $A_{\sigma}\in\mathcal{M}_n(\C)$ dont tous les coefficients sont nuls exceptés ceux d'indice $(\sigma(j),j)$ qui sont égaux à $1$ est diagonalisable. % passer par le morphisme de groupes
			\item Soit $M=\begin{pmatrix}
				A&C\\
				0&B
			\end{pmatrix}$ une matrice triangulaire par blocs. Si $M$ est diagonalisable, alors il en est de même pour $A$ et $B$.
		\end{enumerate}
	\end{ex}
	
	\begin{cor}%cor
		L'endomorphisme $u$ est diagonalisable si, et seulement si, son polynôme minimal est scindé à racines simples.\\
		On a le même résultat pour $A\in\mathcal{M}_n(\K)$.
	\end{cor}
	
	\begin{prop}%prop
		Soit $u\in\mathcal{L}(E)$ diagonalisable et $F$ un sous-espace vectoriel de $E$ stable par $u$. Alors l'endomorphisme induit par $u$ sur $F$ est diagonalisable.
	\end{prop}
	
	\begin{exo}
		Soit $u$ un endomorphisme diagonalisable d'un espace vectoriel $E$ de dimension finie.
		\begin{enumerate}
			\item Montrer que les sous-espaces vectoriels de $E$ stables par $u$ sont les $F=\displaystyle\bigoplus_{\lambda\in\Sp(u)}F_\lambda$, où$F_\lambda$ est un sous-espace vectoriel de $E_\lambda(u)$ pour tout $\lambda\in\Sp(u)$.
			\item A quelle condition y a-t-il un nombre fini de tels sous-espaces vectoriels stables, et dans ce cas, combien y en a-t-il ?
		\end{enumerate}
	\end{exo}
	
	\begin{ex}
		Retrouvons que si la matrice $M=\begin{pmatrix}
			A&C\\
			0&B
		\end{pmatrix}$ est diagonalisable, alors $A$ et $B$ sont diagonalisables.
		\begin{itemize}
			\item La matrice $A\in\mathcal{M}_p(\K)$ représente l'endomorphisme induit, sur le sous-espace vectoriel engendré par les $p$ premiers vecteurs de la base canonique, par l'endomorphisme canoniquement associé à $M$. D'après la proposition précédente, elle est donc diagonalisable.
			\item Pour $B$, ce n'est pas aussi facile car elle n'est pas directement la matrice d'un endomorphisme induit. Mais on peut faire le même raisonnement que précédemment avec la transposée $M^T=\begin{pmatrix}
				A^T&0\\
				C^T&B^T
			\end{pmatrix}$ pour montrer que $B^T$ est diagonalisable.
		\end{itemize}
	\end{ex}
	
	\begin{thm}[Diagonalisation simultanée (HP)]%thm
		Des endomorphismes diagonalisables qui commutent deux à deux sont simultanément diagonalisables.
	\end{thm}
	\begin{proof}
		Récurrence forte sur la dimension de $E$. Distinguer le cas où tous les $u_i$ sont des homothéties. Dans l'autre cas, concaténer les bases obtenues par hypothèse de récurrence pour les sous-espaces propres de $u_{i_0}$ qui n'est pas une homothétie (la commutation fournit l'existence des induits sur ces sous-espaces propres).
	\end{proof}
	
	\begin{exo}
		La réciproque est-elle vraie ?
	\end{exo}
	
	\subsection{Projecteurs spectraux (HP)}
	
	\begin{dfn}
		Soit $u$ un endomorphisme de spectre $\{\lambda_1,\ldots,\lambda_r\}$ où les $\lambda_i$ sont deux à deux distincts. La famille des \textbf{projecteurs spectraux} de $u$ est la famille des projecteurs associés à la décomposition : $$E=\bigoplus_{i=1}^rE_{\lambda_i}(u).$$ On note $(p_{\lambda_1},\ldots,p_{\lambda_r})$ ces projecteurs.
	\end{dfn}
	
	Comme l'endomorphisme induit par $u$ sur ses sous-espaces propres est une homothétie, on a : $$u=\sum_{i=1}^{r}\lambda_ip_{\lambda_i}$$ et, pour tout $n\in\N$ : $$u^n=\sum_{i=1}^{r}\lambda_i^np_{\lambda_i}.$$
	%récurrence ou caractérisation par les supplémentaires ; mieux, vraie pour tout polynôme par linéarité
	
	\begin{ex}
		Puissances de $A=\begin{pmatrix}
			5&9&12\\
			6&8&12\\
			-6&-9&-13
		\end{pmatrix}$.
	\end{ex}
	
	\begin{exo}
		 Si $u$ est diagonalisable et inversible, exprimer $u^n$ en fonction de $p_{\lambda_i}$ pour tout $n\in\Z$.
	\end{exo}
	
	\begin{prop}%prop
		Les projecteurs spectraux de $u$ sont des polynômes en $u$.
	\end{prop}
	\begin{proof}
		Cela peut être vu comme une conséquence du fait (HP) que les projecteurs liés au lemme de décomposition des noyaux sont des polynômes en $u$. Sinon, utiliser l'égalité précédente élargie à tous les polynômes par linéarité, puis considérer la base de Lagrange associée au spectre de $u$.
	\end{proof}
	
	\section{Trigonalisation}
	
	Dans cette section, $u$ est un endomorphisme de $E$, où $E$ est supposé de dimension finie $n$, et $A\in\mathcal{M}_n(\K)$.
	
	\subsection{Définitions}
	
	\begin{dfn}
		\begin{itemize}
			\item L'endomorphisme $u$ est dit \textbf{trigonalisable} s'il existe une base de $E$ dans laquelle la matrice de $u$ est triangulaire supérieure.\\
			Une telle base s'appelle une \textbf{base de trigonalisation} de $u$.
			\item La matrice $A$ est dite \textbf{trigonalisable} si elle est semblable à une matrice triangulaire supérieure, c'est-à-dire s'il existe $T$ triangulaire supérieure et $P$ inversible telles que $A=PTP\inv$.
		\end{itemize}
	\end{dfn}
	
	\paragraph{Interprétation géométrique}

	Dire que la matrice de $u$ dans une base $\mathcal{B}=(e_1,\ldots,e_n)$ est triangulaire supérieure revient à dire que, pour tout $i\in\llbracket0,n\rrbracket$, le sous-espace vectoriel $\Vect(e_1,\ldots,e_n)$ est stable par $u$. Donc l'endomorphisme $u$ est trigonalisable si, et seulement s'il existe des sous-espaces vectoriels $F_0,\ldots,F_n$ de $E$, stables par $u$, tels que : $$\forall i\in\llbracket0,n\rrbracket,\ \dim(F_i)=i\quad\text{et}\quad\forall i\in\llbracket1,n\rrbracket,\ F_{i-1}\subset F_i.$$
	%DRAPEAUX !!!
	
	\begin{terminologie}
		\begin{itemize}
		\item \textbf{Trigonaliser} un endomorphisme $u$ signifie trouver une base dans laquelle sa matrice est triangulaire supérieure.
		\item \textbf{Trigonaliser} une matrice signifie trigonaliser son endomorphisme canoniquement associé, ce qui revient à trouver une matrice $T$ triangulaire supérieure et une matrice $P$ inversible telles que $A=PTP\inv$.
	\end{itemize}
	\end{terminologie}
	
	\begin{ex}
		\begin{enumerate}
			\item Supposons $\K=\R$ et $E$ euclidien. Soit $u$ un endomorphisme trigonalisable de $E$. Considérons une base $\mathcal{B}=(e_1,\ldots,e_n)$ de trigonalisation de $u$. Par le procédé d'orthonormalisation de Gram-Schmidt, on peut alors trouver une base orthonormée $\mathcal{C}=(f_1,\ldots,f_n)$ de $E$ telle que :
			$$\forall i\in\llbracket1,n\rrbracket,\ \Vect(e_1,\ldots,e_i)=\Vect(f_1,\ldots,f_i).$$ D'après l'interprétation ci-dessus, la matrice de $u$ est alors également triangulaire supérieure dans la base $\mathcal{C}$. Autrement dit, un endomorphisme trigonalisable d'un espace euclidien est trigonalisable en base orthonormée.
			\item Soit $u$ un endomorphisme trigonalisable. Sa matrice dans une base $(e_1,\ldots,e_n)$ est donc triangulaire supérieure ; notons $(\lambda_1,\ldots,\lambda_n)$ sa diagonale. Montrons que $\chi_u=(X-\lambda_1)\cdots(X-\lambda_n)$ annule $u$.\\
			Posons $F_k=\Vect(e_1,\ldots,e_k)$ pour tout $k\in\llbracket0,n\rrbracket$, avec donc $F_0=\{0\}$ et $F_n=E$.\\
			Soit $k\in\llbracket1,n\rrbracket$. Par la forme triangulaire supérieure de la matrice de $u$ dans $\mathcal{B}$, on a :
			\begin{itemize}
				\item pour tout $i\in\llbracket1,k-1\rrbracket,\ (u-\lambda_k\Id_E)(e_i)\in F_i\subset F_{k-1}$ ;
				\item $u(e_k)=\lambda_k e_k+x$ avec $x\in F_{k-1}$. Donc $(u-\lambda_k\Id_E)(e_k)\in F_{k-1}$.
			\end{itemize}
			On a alors : $$E=F_n\xrightarrow{u-\lambda_n\Id_E}F_{n-1}\xrightarrow{}\cdots\xrightarrow{}F_1\xrightarrow{u-\lambda_1\Id_E}F_0=\{0\},$$ ce qui donne $((u-\lambda_1\Id_E)\circ\cdots\circ(u-\lambda_n\Id_E))(u)=0$, c'est-à-dire $\chi_u(u)=0$.
		\end{enumerate}
	\end{ex}
	
	\begin{prop}%prop
		Si $A$ représente $u$, elle est trigonalisable si, et seulement si, $u$ est trigonalisable.
	\end{prop}
	
	\begin{cor}%cor
		La matrice $A$ est trigonalisable si, et seulement si, son endomorphisme canoniquement associé est trigonalisable.
	\end{cor}
	
	\begin{rmq}
		On aurait pu aussi choisir la forme triangulaire inférieure. En effet, si la matrice de $u\in\mathcal{L}(E)$ dans la base $(e_1,\ldots,e_n)$e st triangulaire supérieure, alors la matrice de $u$ dans la base $(e_n,e_{n-1},\ldots,e_1)$ est triangulaire inférieure.
	\end{rmq}
	
	\begin{exo}
		Ces deux matrices sont-elles transposées l'une de l'autre ? %non, symétrie centrale
	\end{exo}
	
	\begin{met}
		Une base de trigonalisation commençant par un vecteur propre, pour trigonaliser un endomorphisme $u$, on commence souvent par compléter un vecteur propre de $u$ en une base de $E$, dans laquelle la matrice de $u$ est donc de la forme $$\begin{pmatrix}
			\lambda&L\\
			0&A
		\end{pmatrix}\in\mathcal{M}_n(\K) \ \ \text{où}\ \ A\in\mathcal{M}_{n-1}(\K)\ \ \text{et} \ \ L\in\mathcal{M}_{1,n-1}(\K)$$ Il suffit alors de montrer que l'on peut prendre $A$ triangulaire supérieure pour conclure. Comme la matrice $A$ n'est pas \textit{a priori} la matrice d'un endomorphisme, on raisonne le plus souvent en trigonalisant la matrice $A$ (ou alors on fait le terroriste en composant par une projection sur un hyperplan pour se ramener à un endomorphisme).
	\end{met}
	
	\begin{rmq}
		\begin{itemize}
			\item La démonstration du théorème suivant est une illustration de cette méthode.
			\item Il peut être parfois plus intéressant d'utiliser la forme suivante : 
			$$\begin{pmatrix}
				\lambda I_p&B\\
				0&A
			\end{pmatrix}\in\mathcal{M}_n(\K) \ \ \text{où}\ \ A\in\mathcal{M}_{n-p}(\K)\ \ \text{et} \ \ B\in\mathcal{M}_{p,n-p}(\K)$$ en complétant une base du sous-espace propre $E_\lambda(u)$ en une base de $E$.
		\end{itemize}
	\end{rmq}
	
	\subsection{Utilisation du polynôme caractéristique}
	
	\begin{thm}%thm
		Un endomorphisme $u\in\mathcal{L}(E)$ est trigonalisable si, et seulement si, son polynôme caractéristique est scindé sur $\K$.
	\end{thm}
	\begin{proof}
		Il est équivalent d'établir le résultat pour les matrices ou les endomorphismes.
		\begin{itemize}
			\item Si $A\in\mathcal{M}_n(\K)$ est trigonalisable, alors son polynôme caractéristique est égal à celui d'une matrice triangulaire supérieure : ile st donc scindé sur $\K$.
			\item Montrons la réciproque par récurrence. Posons, pour $n\in\N\st$ :
			\begin{center}
				$H_n$ : "Tout endomorphisme d'un $\K$-espace vectoriel de dimension $n$ dont le polynôme caractéristique est scindé est trigonalisable."\\
				$H_n'$ : "Toute matrice de $\mathcal{M}_n(\K)$ dont le polynôme caractéristique est scindé est trigonalisable."
			\end{center}
			Il est clair que les deux énoncés $H_n$ et $H_n'$ sont équivalents. Ainsi, dans l'hérédité, on démontrera l'implication $H_n'\implies H_{n+1}$. \\
			\textbf{Initialisation :} $H_1$ est immédiat.\\
			\textbf{Hérédité :} Soit $n\geqslant 1$ ; supposons $H_n'$. Soit $E$ un $\K$-espace vectoriel de dimension $n+1$ et $u\in\mathcal{L}(E)$ tel que $\chi_u$ soit scindé. Utilisons le point méthode précédent.\\
			Comme $\chi_u$ est scindé de degré $n+1\geqslant1$, il possède une racine. Par conséquent, il existe $\lambda\in\Sp(u)$. Prenons un vecteur propre associé et complétons-le en une base $\mathcal{B}$ de $E$ de sorte que :
			$$\text{Mat}_\mathcal{B}(u)=\begin{pmatrix}
				\lambda&L\\
				0&A
			\end{pmatrix}\in\mathcal{M}_n(\K) \ \ \text{où}\ \ A\in\mathcal{M}_{n-1}(\K)\ \ \text{et} \ \ L\in\mathcal{M}_{1,n-1}(\K)$$
			Le calcul du polynôme caractéristique de cette matrice triangulaire par blocs donne $\chi_u=(X-\lambda)\chi_A$. Comme $\chi_u$ est scindé sur $\K$, on en déduit que $\chi_A$ l'est également.\\
			D'après $H_n'$, il existe donc $P\in\mathcal{GL}_n(\K)$ tel que $P\inv AP$ soit triangulaire supérieure.\\
			Notons $\begin{pmatrix}
				1&0\\
				0&P
			\end{pmatrix}$. Il s'agit d'une matrice inversible de $\mathcal{M}_{n+1}()\K)$ dont l'inverse est $Q\inv=\begin{pmatrix}
				1&0\\
				0&P\inv
			\end{pmatrix}$ (il suffit de vérifier que le produit vaut $I_{n+1}$).\\
			Un produit par blocs donne alors : $$Q\inv\text{Mat}_\mathcal{B}(u)Q=\begin{pmatrix}
				\lambda&LP\\
				0&P\inv AP
			\end{pmatrix}.$$ Ainsi, $Q\inv\text{Mat}_\mathcal{B}(u)Q$, qui est la matrice de $u$ dans une certaine base, est triangulaire supérieure, ce qui permet de conclure que $u$ est trigonalisable. Cela prouve $H_{n+1}$.
		\end{itemize}
	\end{proof}
	
	\begin{rmq}
		De plus, on voit, en adaptant la démonstration ci-dessous, que quel que soit l'ordre dans lequel on écrit la liste $(\lambda_1,\ldots,\lambda_n)$ des valeurs propres de $u$, comptées avec ordre de multiplicité, on peut trouver une base de $E$ dans laquelle la matrice de $u$ est triangulaire supérieure avec $(\lambda_1,\ldots,\lambda_n)$ pour diagonale.
	\end{rmq}
	
	\begin{ex}
		\begin{enumerate}
			\item Comme une matrice et sa transposée ont même polynôme caractéristique, on en déduit qu'une matrice $A$ est trigonalisable, si, et seulement si, $A^T$ est trigonalisable.
			\item Si $u$ est trigonalisable et si $F$ est un sous-espace vectoriel de $E$ stable par $u$, alors l'endomorphisme induit $u_F$ est aussi trigonalisable.\\
			En effet, le polynôme caractéristique de $u_F$ divise celui de $u$ qui est scindé. Ainsi, le polynôme caractéristique de $u_F$ est scindé donc $u_F$ est trigonalisable.
		\end{enumerate}
	\end{ex}
	
	\begin{cor}%cor
		Si $E$ est un $\C$-espace vectoriel de dimension finie non nulle, alors tout endomorphisme de $E$ est trigonalisable. \\
		Toute matrice carrée à coefficients dans $\C$ est trigonalisable sur $\C$.
	\end{cor}
	\begin{proof}
		Le théorème de d'Alembert-Gauss nous dit que le polynôme caractéristique est scindé, ce qui permet de conclure d'après le théorème précédent.
	\end{proof}
	
	\begin{rmq}
		On peut en déduire le théorème de Cayley-Hamilton lorsque $\K$ est un sous-corps de $\C$. En effet, un exemple précédent montre que le polynôme caractéristique est un polynôme annulateur pour un endomorphisme trigonalisable, donc pour tout endomorphisme d'un $\C$-espace vectoriel de dimension finie.C'est alors vrai pour toute matrice de $\mathcal{M}_n(\C)$, et donc également pour toute matrice de $\mathcal{M}_n(\K)$ puis pour tout endomorphisme d'un $\K$-espace vectoriel de dimension finie.
	\end{rmq}
	
	\begin{ex}
		Trigonalisons $A=\begin{pmatrix}
			0&3&3\\
			-1&8&6\\
			2&-14&-10
		\end{pmatrix}$.
	\end{ex}
	
	"Cette histoire de compléter en une base, ça me fait penser à un exo de ma petite fille qui est en CP."
	
	\begin{exo}
		\begin{enumerate}
			\item La matrice $A=\begin{pmatrix}
				14&18&18\\
				-6&-7&-9\\
				-2&-2&-1
			\end{pmatrix}\in\mathcal{M}_3(\R)$ est-elle diagonalisable ? trigonalisable ?
			\item Montrer que la matrice $A$ est semblable à la matrice $T=\begin{pmatrix}
				2&0&0\\
				0&2&1\\
				0&0&2
			\end{pmatrix}$ et expliciter une matrice inversible $P$ telle que $A=PTP\inv$.
		\end{enumerate}
	\end{exo}
	
	\subsection{Utilisation d'un polynôme annulateur}
	
	\begin{rmq}
		Rappelons qu'un polynôme scindé est \textit{non nul} par définition.
	\end{rmq}
	
	\begin{thm}%thm
		L'endomorphisme $u$ est trigonalisable si, et seulement s'il possède un polynôme annulateur scindé.\\
		On a le même résultat pour $A\in\mathcal{M}_n(\K)$.
	\end{thm}
	\begin{proof}
		\begin{itemize}
			\item Si $u$ est trigonalisable, son polynôme caractéristique est scindé, et il est annulateur d'après le théorème de Cayley-Hamilton.
			\item Pour la réciproque, on procède par récurrence à l'aide du point méthode sur la trigonalisation.
		\end{itemize}
	\end{proof}
	
	\begin{cor}%cor
		L'endomorphisme $u$ est trigonalisable si, et seulement si, son polynôme minimal est scindé.\\
		On a le même résultat pour $A\in\mathcal{M}_n(\K)$.
	\end{cor}
	
	\begin{thm}[Trigonalisation simultanée (HP)]%thm
		Des endomorphismes trigonalisables qui commutent deux à deux sont simultanément trigonalisables.
	\end{thm}
	\begin{proof}
		Par récurrence forte sur la dimension de $E$. Le résultat est équivalent à sa version matricielle, on démontre l'implication des endomorphismes vers les matrices dans l'hérédité.
	\end{proof}
	
	\subsubsection*{Cas des endomorphismes nilpotents}
	
	\begin{prop}{}%prop
		\begin{itemize}
			\item Un endomorphisme est nilpotent si, e seulement s'il est trigonalisable et toutes ses valeurs propres sont nulles.
			\item Une matrice est nilpotente si, et seulement si, elle est semblable à une matrice triangulaire supérieure stricte.
		\end{itemize}
	\end{prop}
	\begin{proof}
		Le deuxième point est simplement la traduction matricielle du premier.
		\begin{itemize}
			\item Un endomorphisme nilpotent est annulé par un polynôme scindé de la forme $X^k$, donc est trigonalisable. D'autre part, ses valeurs propres sont toutes racines de $X^k$, donc sont toutes nulles.
			\item Soit $u$ un endomorphisme trigonalisable ayant toutes ses valeurs propres nulles. Sont polynôme caractéristique est donc $X^n$ et d'après le théorème de Cayley-Hamilton, on a donc $o^n=0$, ce qui prouve que $u$ est nilpotent.
		\end{itemize}
	\end{proof}
	
	\begin{rmq}
		En particulier, une matrice triangulaire supérieure stricte est nilpotente.
	\end{rmq}
	
	\begin{cor}{}%cor
		\begin{itemize}
			\item Si $\dim(E)=n$, un endomorphisme de $E$ est nilpotent si, et seulement si, son polynôme caractéristique est $X^n$.
			\item Une matrice $A\in\mathcal{M}_n(\K)$ est nilpotente si, et seulement si, son polynôme caractéristique est $X^n$.
		\end{itemize}
	\end{cor}
	
	\subsection{Sous-espaces caractéristiques}
	
	\begin{prop}[Décomposition en sous-espaces caractéristiques]
		Soit $u\in\mathcal{L}(E)$ dont le polynôme caractéristique $\chi_u=\displaystyle\prod_{\lambda\in\Sp(u)}(X-\lambda)^{m(\lambda)}$ est scindé. On a la décomposition : $$E=\bigoplus_{\lambda\in\Sp(u)}\ker(u-\lambda\Id_E)^{m(\lambda)}.$$
	\end{prop}
	\begin{proof}
		D'après le théorème de Cayley-Hamilton, on a $\chi_u(u)=0$ et le lemme de décomposition des noyaux donne alors le résultat puisque les polynômes $(X-\lambda)^{m(\lambda)}$ sont premiers entre eux deux à deux.
	\end{proof}
	
	\begin{terminologie}
		Les sous-espaces vectoriels $F_\lambda=\ker(u-\lambda\Id_E)^{m(\lambda)}$, pour $\lambda\in\Sp(u)$, sont appelés \textbf{sous-espaces caractéristiques} de l'endomorphisme de $u$.
	\end{terminologie}
	
	Soit $\lambda\in\Sp(u)$. Comme noyau d'un polynôme en $u$, $F_\lambda$ est stable par $u$ et l'endomorphisme induit $u_\lambda$ par $u$ sur $F_\lambda$ vérifie $(u_\lambda-\lambda\Id_{F_\lambda})^{m(\lambda)}=0$. Ainsi, $u_\lambda$ est la somme d'une homothétie de rapport $\lambda$ et d'un endomorphisme nilpotent.
	
	\begin{rmq}
		Soit $\pi_u$ le polynôme minimal d'un tel endomorphisme. Il divise $chi_u$ donc de la forme $\pi_u=\displaystyle\prod_{\lambda\in\Sp(u)}(X-\lambda)^{n(\lambda)}$, avec $n(\lambda)\leqslant m(\lambda)$ pour tout $\lambda\in\Sp(u)$. On a donc également : $$E=\bigoplus_{\lambda\in\Sp(u)}\ker(u-\lambda\Id_E)^{n(\lambda)}.$$
	\end{rmq}
	
	\begin{prop}{}%prop
		Pour toute valeur propre $\lambda$, la dimension de $F_\lambda$ est l'ordre de multiplicité de $\lambda$.
	\end{prop}
	
	\begin{cor}{}%cor
		Soit $u\in\mathcal{L}(E)$ dont le polynôme caractéristique est scindé. Il existe une base de $E$ dans laquelle la matrice de $E$ est diagonale par blocs, chaque bloc étant triangulaire à coefficients diagonaux égaux.
	\end{cor}
	\begin{proof}
		Dans chaque sous-espace caractéristique $F_\lambda$, on choisit une base de trigonalisation de l'endomorphisme nilpotent induit par $u-\lambda\Id_E$ sur $F_\lambda$.
	\end{proof}
	
	\begin{ex}
		Soit $u$ l'endomorphisme de $\R^3$ canoniquement associé à la matrice : $$A=\begin{pmatrix}
			0&-3&-5\\
			-1&0&-3\\
			1&2&5
		\end{pmatrix}\in\mathcal{M}_3(\R).$$
		On a $$A=PTP\inv\quad\text{avec}\quad T=\begin{pmatrix}
			1&0&0\\
			0&2&1\\
			0&0&2
		\end{pmatrix}\quad\text{et}\quad P=\begin{pmatrix}
		-2&-1&-1\\
		-1&-1&1\\
		1&1&0
		\end{pmatrix}.$$
	\end{ex}
	
	\subsubsection*{Décomposition de Dunford (HP)}
		Supposons $\chi_u$ scindé. Le corollaire précédent fournit une base $\mathcal{B}$ de $E$ dans laquelle la matrice de $u$ est $A=\text{Diag}(T_1,\ldots,T_r)$, où chaque matrice $T_i$ est somme d'une matrice scalaire $D_i$ et d'une matrice triangulaire supérieure stricte $N_i$.
		
		En prenant $D=\text{Diag}(D_1,\ldots,D_r)$ et $N=\text{Diag}(N_1,\ldots,N_r)$, on obtient $A=D+N$, avec $D$ diagonale, et $N$ triangulaire supérieure stricte, donc nilpotente.
		
		De plus, pour tout $i\in\llbracket1,r\rrbracket$, on a $D_iN_i=N_iD_i$, puisque $D_i$ est scalaire. Donc, par produit de matrices diagonales par blocs, cela donne $DN=ND$.
		
		
		En notant $d$ (respectivement $n$) l'endomorphisme de $E$ dont la matrice dans $\mathcal{B}$ (respectivement $N$), on obtient $u=d+n$, avec $d$ diagonalisable et $n$ nilpotent commutant avec $d$.
		
		
		L'exercice Moulin.Red.63 montre l'unicité d'un tel couple $(d,n)$. Cette décomposition $u=d+n$ s'appelle la \textbf{décomposition de Dunford} de $u$. On donne ci-dessous une autre démonstration de l'unicité de la décomposition de Dunford, qui passe par le théorème de diagonalisation simultanée.
	
	\begin{ex}[Unicité de la décomposition de Dunford par la codiagonalisation]
		On a trouvé $d_0$ diagonalisable et $n_0$ nilpotent qui commutent et tels que $u=d_0+n_0$.\\
		Or, on a : $$d_0=\sum_{\lambda\in\Sp(u)}\lambda p_\lambda$$ donc $d_0\in\K[u]$ puis $d_0$ commute avec $u$, et donc $n_0$ aussi.\\
		Soit $(d,n)$ une décomposition de Dunford de $u$. Alors $d$ et $n$ commutent avec $u$, donc avec $d_0$ et $n_0$. \\
		Donc, d'après le théorème de diagonalisation simultanée, $d$ et $d_0$ sont simultanément diagonalisables, donc $d-d_0$ est diagonalisable. Mais $n_0-n$ est nilpotent comme différence de deux endomorphismes nilpotents qui commutent.\\
		Ainsi, comme $u=d_0+n_0=d+n$, on a montré que $d-d_0=n_0-n$ est nilpotent et diagonalisable, donc nul.\\
		Ainsi, $d=d_0$ et $n=n_0$, si bien que la décomposition de Dunford est unique.
		
	\end{ex}
	
\end{document}